\begin{header}
    \fontsize{25pt}{25pt}\selectfont \ResumeName

    \vspace{5pt}

    \normalsize
    \ifdefined\ResumeEmail
        \mbox{\hrefWithoutArrow{mailto:\ResumeEmail}{\ResumeEmail}}\\
    \fi
    \mbox{\ResumeLocation}%
    \ifdefined\ResumePhone
        \ResumeContactSep\mbox{\ResumePhone}%
    \fi
    \ResumeContactSep\mbox{\hrefWithoutArrow{https://\ResumeLinkedIn}{\ResumeLinkedIn}}%
    \ResumeContactSep\mbox{\hrefWithoutArrow{https://\ResumeGitHub}{\ResumeGitHub}}%
\end{header}

\vspace{4pt}

\section{Summary}
Aerospace engineer with a PhD in Aerospace Engineering from Georgia Tech, focused on computational modeling, structural and thermal finite-element analysis, and multidisciplinary design optimization. Experience at NASA Glenn Research Center includes structural and thermal analysis, battery pack thermal management, hybrid-electric propulsion modeling, and modernization of NASA’s Chemical Equilibrium with Applications software. Strong background in developing validated engineering tools, high-performance analysis workflows, and open-source research software.

\section{Education}
\begin{threecolentry}{\textbf{Georgia Institute of Technology}}{August 2021}
    PhD, Aerospace Engineering
    \begin{highlights}
        \item Thesis: \href{http://hdl.handle.net/1853/65093}{\textit{Stress-Based Topology Optimization for Steady-State and Transient Thermoelastic Design}}
    \end{highlights}
\end{threecolentry}

\begin{threecolentry}{}{May 2018}
    MS, Aerospace Engineering
\end{threecolentry}

\begin{threecolentry}{\textbf{The University of Texas at Austin}}{May 2016}
    BS, Aerospace Engineering
\end{threecolentry}

\section{Experience}
\textbf{NASA Glenn Research Center}
\begin{twocolentry}{June 2022 -- Present}
    Aerospace Engineer -- Cleveland, OH
\end{twocolentry}

\underline{Battery Pack Thermal Analysis and Hybrid-Electric Propulsion:}
\begin{onecolentry}
    \begin{highlights}
        \item Developed a hybrid-electric propulsion analysis model coupling time-varying battery cell-discharge and electric motor simulations, enabling aircraft-level optimization studies for reduced climb fuel consumption.
        \item Implemented phase-change material constitutive modeling in \href{https://github.com/smdogroup/tacs}{TACS} for transient battery pack finite-element simulations, enabling evaluation of passive thermal management strategies.
        \item Developed a two-way, time-resolved coupling framework for battery pack shape optimization, integrating transient finite-element heat transfer with cell-discharge models.
    \end{highlights}
\end{onecolentry}

\underline{Modernization of Chemical Equilibrium with Applications \href{https://github.com/nasa/cea}{CEA}:}
\begin{onecolentry}
    \begin{highlights}
        \item Led modernization of NASA’s Chemical Equilibrium with Applications software, refactoring a legacy Fortran codebase into a modern, maintainable scientific software package for rocket, shock, and detonation analysis.
        \item Led public open-source release preparation for NASA CEA, including validation workflows, interface documentation, packaging improvements, and repository maintenance.
        \item Maintain NASA's open-source CEA repository, managing public releases, user-facing interfaces, validation updates, and long-term software modernization.
        \item Developed C, Python, and MATLAB interfaces to the Fortran computational core, enabling CEA integration with modern analysis, optimization, and scripting workflows.
        \item Optimized the Python interface to execute more than 100,000 equilibrium calculations in under 2 seconds, supporting high-throughput design and optimization workflows.
        \item Developed analytic derivative capabilities for equilibrium thermodynamic outputs, enabling gradient-based integration of CEA into multidisciplinary design optimization workflows.
    \end{highlights}
\end{onecolentry}

\vspace{0.2cm}

\textbf{Peerless Technologies}
\begin{twocolentry}{June 2021 -- June 2022}
    Aerospace Engineer (Contractor at NASA GRC) -- Cleveland, OH
\end{twocolentry}

\vspace{0.10cm}
\begin{onecolentry}
    \begin{highlights}
        \item Developed propeller optimization capabilities using GXBeam.jl and automatic differentiation, enabling coupled structural and acoustic analysis of design concepts.
        \item Conducted transient finite-element heat-conduction analyses to predict battery temperatures during thermal runaway, achieving strong agreement with experimental test data.
        \item Implemented transient heat-conduction analysis capabilities in the open-source \href{https://github.com/OpenMDAO/mphys}{MPhys} framework using the \href{https://github.com/smdogroup/tacs}{TACS} finite-element solver.
    \end{highlights}
\end{onecolentry}

\section{Selected Publications}
\begin{onecolentry}
    \begin{highlights}
      \item Psenica, C., Fang, L., Zoppelt, S., \textbf{Leader, M.K.}, and He, P., A Modular Conjugate Heat Transfer Optimization Framework for Thermal Management of Electric Aircraft. \textit{International Journal of Heat and Mass Transfer}, 259:128345, 2026.
      \item \textbf{Leader, M.K.}, Battery Pack Shape Optimization Using Transient Heat Conduction Coupled with Cell-Discharge Analysis. \textit{AIAA SCITECH 2025 Forum}, Orlando, FL, 2025. doi:\href{https://arc.aiaa.org/doi/abs/10.2514/6.2025-0361}{10.2514/6.2025-0361}. AIAA 2025-0361.
      \item \textbf{Leader, M.K.}, Aretskin-Hariton, E., and Moore, K.T., Multidisciplinary Optimization of a Transonic Truss Braced Wing Aircraft with Hybrid-Electric Propulsion. \textit{AIAA AVIATION Forum and ASCEND 2024}, Las Vegas, NV, 2024. doi:\href{https://arc.aiaa.org/doi/abs/10.2514/6.2024-4294}{10.2514/6.2024-4294}. AIAA 2024-4294.
      \item \textbf{Leader, M.K.}, Lavelle, T.M., Wang, X.J., Dickens, K.W., McTague, M., and Hill, J.P., CEA2022: A Modernization of NASA Glenn's Software CEA (Chemical Equilibrium with Applications). \textit{Thermal and Fluids Analysis Workshop (TFAWS) 2024}, Cleveland, OH, 2024. \href{https://ntrs.nasa.gov/citations/20240009728}{NASA NTRS}.
      \item \textbf{Leader, M.K.} and Kennedy, G.J., Thermoelastic Topology Optimization Using Steady-State and Transient Analysis for Stress and Thermal Performance. \textit{AIAA SciTech 2021 Forum}, Virtual, 2021. doi:\href{https://arc.aiaa.org/doi/abs/10.2514/6.2021-1895}{10.2514/6.2021-1895}. AIAA 2021-1895.
      \item \textbf{Leader, M.K.}, Chin, T.W., and Kennedy, G.J., High-Resolution Topology Optimization with Stress and Natural Frequency Constraints. \textit{AIAA Journal}, 2019. doi:\href{https://arc.aiaa.org/doi/10.2514/1.J057777}{10.2514/1.J057777}.
      \item Chin, T.W., \textbf{Leader, M.K.}, and Kennedy, G.J., A Scalable Framework for Large-Scale 3D Multimaterial Topology Optimization with Octree-Based Mesh Adaptation. \textit{Advances in Engineering Software}, 2019. doi:\href{https://www.sciencedirect.com/science/article/abs/pii/S0965997818309682}{10.1016/j.advengsoft.2019.05.004}.
    \end{highlights}
\end{onecolentry}

\section{Skills}
\begin{onecolentry}
    \textbf{Methods:} Finite-element analysis (FEA), transient heat transfer, structural dynamics, buckling analysis, topology optimization, shape optimization, multidisciplinary design optimization (MDAO), chemical equilibrium analysis, high-performance computing (HPC)
\end{onecolentry}

\begin{onecolentry}
    \textbf{Programming:} Python, Fortran, C, C++, MATLAB, Julia, Bash, Perl, Cython
\end{onecolentry}

\begin{onecolentry}
    \textbf{Software and frameworks:} \href{https://github.com/OpenMDAO/OpenMDAO}{OpenMDAO}, \href{https://github.com/smdogroup/tacs}{TACS}, git, MPI, Nastran, Femap, SolidWorks, NX CAD
\end{onecolentry}
